% Chapter 6 - Educational and Didactic Concept
% Ehemals Kapitel 5, jetzt nach CrypTool Online Architecture (Rueckmeldung
% des Betreuers).
\section{Educational and Didactic Concept}
\label{chap:educationalConcept}

\subsection{Learning Objectives}
\label{sec:learningObjectives}
% TODO: siehe Strukturplan.tex, Abschnitt 6.1
% - PQC verstehen, McEliece einordnen koennen, Schluesselablaeufe nachvollziehen koennen

\subsection{Target Audience Analysis}
\label{sec:targetAudience}
% TODO: siehe Strukturplan.tex, Abschnitt 6.2
% - Bachelor-Studierende, Grundkenntnisse Kryptographie

\subsection{Educational Challenges}
\label{sec:educationalChallenges}
% TODO: siehe Strukturplan.tex, Abschnitt 6.3
% - Abstrakte Mathematik, hohe Komplexitaet

\subsection{Visualization Principles}
\label{sec:visualizationPrinciples}
% TODO: siehe Strukturplan.tex, Abschnitt 6.4
% - Cognitive Load Theory, Progressive Disclosure, Step-by-Step Visualization

\subsection{Concept of the Learning Tool}
\label{sec:learningToolConcept}
% TODO: siehe Strukturplan.tex, Abschnitt 6.5
% - Lernszenarien, Benutzerfluss, Interaktionen
